\section{결 론}

본 논문에서는 실제 데이팅 앱의 매칭 데이터를 활용하여 시각적 매칭 호환성을 예측하는 딥러닝 기반 시스템을 제안하였다. Qwen3-VL 백본의 파라미터를 동결하고 경량 프로젝션 헤드만 학습하는 효율적인 구조를 설계하였으며, InfoNCE Loss 기반 대조 학습을 통해 실제 커플 쌍을 임베딩 공간에서 가깝게 배치하도록 학습하였다.

실험 결과, 제안 모델은 사전 학습된 VLM 베이스라인(R@1 0.81\%) 대비 109\% 향상된 R@1 1.69\%를 달성하였으며, 이는 랜덤 기준 대비 2.6배에 해당한다. 본 연구를 통해 다음의 인사이트를 도출하였다. 첫째, 사전 학습된 VLM의 일반적 특징만으로는 커플 매칭 예측에 한계가 있으며, 도메인 특화 프로젝션 헤드 학습이 필수적이다. 둘째, 소규모 데이터셋에서는 Dropout이 오히려 역효과를 보인다. 셋째, InfoNCE Loss의 Temperature 완화($\tau$=0.07$\rightarrow$0.1)와 Weight Decay 증가가 일반화 성능 향상에 효과적이다.

향후 연구로는 텍스트 임베딩(성격, 가치관)과 이미지 임베딩을 결합한 하이브리드 모델 개발, 그리고 더 많은 커플 데이터 확보를 통한 일반화 성능 향상을 계획한다.
